
\chapter*{Comment travailler avec ce cours}
\addcontentsline{toc}{chapter}{Comment travailler avec ce cours}
Ce volume concerne la \textbf{première générale, spécialité mathématiques},
selon le programme applicable à la rentrée 2026. Il prolonge le cours de seconde
de la collection et prépare un parcours jusqu'au master, en passant par la
terminale et la licence. Il ne prétend pas enseigner dès la première tout le
programme d'un master.

Chaque chapitre donne ses objectifs et prérequis, ses définitions, des propriétés
démontrées et des exemples, puis explique l'utilité actuelle et future des outils.
La progression n'est pas un calendrier rigide : probabilités, géométrie, calcul
et raisonnement peuvent être alternés pendant l'année.
Le guide contient la table de couverture officielle, la progression proposée,
les automatismes, les sources et les dépendances des démonstrations.

\section*{Trois passages conseillés}
\begin{enumerate}
\item Lire le socle, refaire les exemples sans regarder, puis résoudre les premiers
exercices de chaque chapitre. Les chapitres 1 à 14 et 16 à 24 portent les thèmes
du programme, avec des compléments explicitement signalés.
\item Reconstruire les démonstrations : annoncer les hypothèses, expliquer chaque
implication, vérifier le domaine et les cas limites. Comparer seulement ensuite
avec le corrigé. Les six exercices de chaque chapitre ont la même numérotation
dans le recueil et dans les corrigés.
\item Explorer le chapitre 15 et les chapitres 25 à 30, hors programme.
Ils sont conçus pour être relus après la terminale. Ils construisent des liens
vers l'analyse, l'algèbre linéaire, les probabilités et le calcul scientifique.
\end{enumerate}

\section*{Le statut des énoncés}
Une définition fixe le sens d'un mot ; elle ne se démontre pas.
Une propriété exige une preuve à partir d'hypothèses.
Un exemple, une figure ou un programme ne suffisent pas à établir un résultat
pour tous les réels ou tous les entiers.
Les fondements admis sont indiqués : propriétés usuelles des nombres,
complétude de $\R$, principe de récurrence et acquis géométriques.
L'existence de l'exponentielle, admise dans le socle officiel, est démontrée
au chapitre 26. Les bornes atteintes utilisées pour Rolle au chapitre 10
sont établies au chapitre 25. Ces renvois évitent de masquer une dépendance.

Les symboles $\forall$ et $\exists$, les limites formelles, les séries et les
espaces vectoriels sont des approfondissements, non des exigences supplémentaires
attribuées au programme officiel. Les programmes exécutables se trouvent dans
\texttt{premiere/python/} ; ils complètent les démonstrations sans les remplacer.
